Істина не десь там — хоч би де було те «там». Істина не міститься ні в релігійних обрядах, ні в таємних доктринах, ні в дотику гуру чи його блаженній усмішці, ні в екзотичних краях чи давніх храмах. Істина — це буквально єдине, що справді існує. Вона не схована, а лежить на видноті; її не бракує — вона присутня в надлишку.

Абсолютна Істина — не вірування, не релігія, не філософія, не миттєвий досвід і не минуще духовне переживання. Вона не статична і не в русі, не добра і не зла. Вона інша, ніж усе це, — інша настільки, що цього не уявити. Думка не може її торкнутися, розум не може її виобразити. Її можна знайти лише в серцевині вселенського буття. Пізнати себе — ось ключ. Виявити своє буття — ось Шлях.

У вченні «Шлях звільнення» є три основні практики, які діють разом із п'ятьма основами й допомагають виявити й здійснити позачасову Істину. Ці основні практики можуть здаватися простими, але не обманюйтеся: поєднані й застосовані всім серцем, вони можуть бути надзвичайно потужними.

Наш розум може вважати, що для проводу до Дійсності нам потрібні тонкі й складні духовні вчення, але це не так. Насправді що складніше вчення, то легше розуму сховатися від себе самого в цій складності, уявляючи, ніби він просувається до просвітлення. Але часто він просувається лише в тому, що витворює дедалі заплутаніші кола, якими ходить і ходить навколо.

Незамінний елемент будь-якого духовного вчення лежить не у вченні, а в щирості й безстрашності людини, яка його застосовує. Навіть якщо часом ви почуваєтеся цілком загубленими у власній дурноті, як сказав Вільям Блейк: «Дурень, що вперто тримається своєї дурноти, стане мудрецем».

Думайте про духовну практику як про своєрідну «прикладну дурноту».

Основні практики — це те, до чого треба відчути смак, приблизно так, як учаться відчувати рівновагу, коли вчаться їздити на велосипеді. Щоб вони справді діяли, вони мають стати частиною вас. Постава, з якою ви їх застосовуєте, не менш важлива за саму практику — тобто вам треба знайти такий спосіб застосування основних практик, який пасує вашому темпераменту й особистому стилю. Ніхто не скаже вам точно, як це робити. Ви просто відкриваєте це спробами й помилками. І те, як ви використовуєте основні практики, розвиватиметься разом із вашим рівнем здійснення.

Не застосовуйте основні практики надто свавільно й не докладайте величезних зусиль. Їх слід застосовувати молитовно, тобто з великою щирістю й відкритістю розуму та серця. Часом вас дуже викликатиме те, що відкривають ці вчення, а часом ви боротиметеся з власною сум'яттю й сумнівом — але пам'ятайте, що благодать надзвичайно важлива й присутня завжди. І часто найтемніше буває перед світанком.

Усі наші найбільші прориви, здається, приходять несподівано, коли ми чекаємо їх найменше. Раптом нам дарують — хмари сум'яття розходяться, і ми бачимо з незвичною ясністю й свободою. Така благодать ніколи не відкладена, її не заслуговують і не виборюють. Її не дають одним і не дають іншим. Благодать присутня завжди; змінюється лише наша відкритість до неї. Певною мірою «Шлях звільнення» — це засіб відкритися благодаті.

Три основні практики — це медитація, самодослідження і споглядання.

\section*{Медитація}

У різних формах езотеричної духовності медитацію часто або перебільшують, або применшують. Коли медитацію перебільшують як єдиний шлях до просвітлення, увага надто часто зосереджується на спробах досягти певного медитативного стану. Остання Дійсність — не якийсь стан свідомості, хоч би яким чудовим чи блаженним він був. Дійсність — це основа всього буття, непороджена й невмируща вічність. Вона присутня в одному досвіді чи стані свідомості так само, як і в будь-якому іншому. Дійсність, або Істина, — це те, що остаточно істинне в усіх станах, у всі часи, в усіх місцях.

На іншому краї — учення, які применшують цінність медитації. Логіка така: оскільки Дійсність присутня завжди й усюди, то медитацією нічого не досягти. Більше того, мовляв, медитація лише підсилюватиме вірування, ніби людина відокремлена від Дійсності й мусить щось робити, щоб її досягти. У цій думці є логіка, але вона може привести до фаталізму й суто інтелектуального розуміння, що перешкоджає справжньому пробудженню. Хоча немає жодного шляху чи практики, які ведуть прямо до пробудження, так само правда й те, що те, що ви робите, життєво важливе для перебігу вашого духовного життя. Рівновага — ось ключ, зусилля без зусиль — ось Шлях.

Медитація — не засіб досягти мети і не щось таке, що можна вдосконалити. Правильно виконувана медитація — це вияв Дійсності, а не шлях до неї. Медитація, виконувана неправильно, — це досконале дзеркало того, як ви опираєтеся теперішній миті, оцінюєте її чи прив'язуєтеся до неї. Медитація діє як досконале дзеркало, що відбиває ваші стосунки із собою, з життям і з теперішньою миттю. Якщо ви близько й уважно побачите, як опираєтеся вмісту теперішньої миті або прив'язуєтеся до нього і наскільки це марно, ви можете відкрити, що означає справді відпустити весь свій опір теперішній миті.

У «Шляху звільнення» медитація має цілком конкретне визначення, призначення й застосування. Медитація — це мистецтво дозволяти всьому просто бути якнайглибше. Щоб дозволити всьому бути, ми мусимо відпустити зусилля контролювати й керувати своїм досвідом — а це означає відпустити особисту волю. Це б'є просто в серце егоїстичної будови, яка шукає щастя через контроль, пошук, намагання й маніпуляцію. Багато форм медитації ґрунтуються на тому, що людина вчиться контролювати свій досвід як засіб досягти спокою. Такі методи часто заводять у безвихідь, де спокій розуму триває лише доти, доки его стримує медитативна техніка.

Тиша й нерухомість медитації — це наріжний камінь, на якому спочиває це вчення. Вона плекає внутрішню стійкість, об'єктивність, неприв'язаність і глибину розуміння, невідому понятійному розуму. Формальну медитацію найкраще робити сидячи (або лежачи, якщо це конче потрібно) там, де вас не потурбують. Постава, що сприяє медитації, — це віддання, відсутність зусиль і відкритість.

Медитація радше форма мовчазної молитви, ніж техніка, яку треба опанувати. Починаючи медитувати, ви можете присвятити десять-п'ятнадцять хвилин тому, щоб просто бути в тиші. Коли цей час стане зручним, можете подовжувати медитацію п'ятихвилинними кроками, аж доки зможете спокійно сидіти тридцять чи сорок хвилин. Але навіть п'ятнадцять-двадцять хвилин тиші щодня почнуть формувати у вас внутрішню зібраність, нерухомість і стійкість.

Важливо розуміти, що, медитуючи, ви даєте зобов'язання чомусь іншому, ніж ваш невгамовний розум. Медитація — не час розбиратися в речах чи аналізувати свій досвід. Не слід і боротися з розумом або намагатися його заспокоїти. Просто спостерігайте думки так, як спостерігали б хмари, що пливуть небом. У ваших думках немає нічого особистого. Це лише явища, що проходять крізь усвідомлення. Медитація — не техніка, яку треба опанувати; це найвища форма молитви, голий акт любові й віддання без зусиль у мовчазну безодню за межами всякого знання.

Медитація не обмежена часами формального сидіння; найглибше вона — це постава буття, спочивання в бутті й як буття. Коли ви відчуєте її смак, то зможете налаштовуватися на неї дедалі частіше в повсякденному житті. Зрештою, у стані визволення медитація просто стане вашим природним станом.

Нижче — опис того, що я називаю Справжньою медитацією. Прочитайте його й дозвольте йому відкрити вам своє справжнє значення, втілюючи його в практику. З часом ви дедалі глибше розумітимете, що таке Справжня медитація. Можливо, вам захочеться щоразу перечитувати «Справжню медитацію» безпосередньо перед тим, як сісти медитувати, доки ви не відчуєте, що вказівки стали вашими.

\subsection*{Справжня медитація}

Справжня медитація не має напрямку й мети. Це чисте віддання без слів, чиста мовчазна молитва. Усі методи, скеровані на досягнення певного стану розуму, обмежені, минущі й зумовлені. Захоплення станами веде лише до неволі й залежності. \emph{Справжня медитація — це нерухомість без зусиль, перебування первісним буттям.}

Справжня медитація з'являється у свідомості спонтанно, коли усвідомленням не маніпулюють і його не контролюють. Коли ви тільки починаєте медитувати, то помічаєте, що увагу часто тримають у полоні, зосереджуючись на якомусь об'єкті: на думках, тілесних відчуттях, емоціях, спогадах, звуках тощо. Так відбувається тому, що розум зумовлений зосереджуватися на об'єктах і стискатися навколо них. Далі розум нав'язливо тлумачить те, що усвідомлює (об'єкт), і намагається ним керувати — механічно й викривлено. Він починає робити висновки й припущення згідно з минулим зумовленням.

У Справжній медитації всі об'єкти (думки, почуття, емоції, спогади тощо) залишені їхньому природному функціонуванню. Це означає, що не слід докладати жодних зусиль, щоб зосереджуватися на об'єкті усвідомлення, маніпулювати ним, контролювати його чи придушувати. У Справжній медитації наголос на тому, щоб \emph{бути} усвідомленням — не на тому, щоб усвідомлювати об'єкти, а на тому, щоб спочивати як сама свідома присутність. Медитуючи, ви не намагаєтеся змінити свій досвід; ви змінюєте свої стосунки з досвідом.

Коли ви м'яко розслабляєтеся в усвідомленні, нав'язливе стискання розуму навколо об'єктів згасає. Тиша буття дедалі ясніше приходить у свідомість як запрошення спочити й перебувати. Постава відкритої сприйнятливості, вільна від будь-якої мети чи передчуття, допоможе тиші й нерухомості відкритися як вашому природному станові.

Коли ви без зусиль дедалі глибше спочиваєте в нерухомості, усвідомлення звільняється від нав'язливої звички розуму контролювати, стискатися й ототожнюватися. Усвідомлення повертається до свого природного стану свідомого буття, абсолютного непроявленого потенціалу — мовчазної безодні за межами всякого знання.

\subsection*{Поширені запитання про медитацію}

\textbf{З.} Мені здається, що центральна вказівка Справжньої медитації — просто перебувати мовчазним, нерухомим усвідомленням. Однак я часто помічаю, що застрягаю в розумі. Чи можна використовувати більш скеровану медитацію, наприклад стеження за диханням, щоб мати на чому зосередитися й не губитися в розумі?

\textbf{В.} Цілком можна використовувати скерованішу техніку — скажімо, стеження за диханням, просту мантру чи молитву-осередження, — якщо ви відчуваєте, що це допомагає не губитися в думках. Але завжди схиляйтеся до дедалі меншої кількості технік. Знаходьте час у кожному медитативному періоді, щоб просто спочивати мовчазним, нерухомим усвідомленням. Справжня медитація — це поступове відпускання того, хто медитує, без того, щоб губитися в розумі.

\textbf{З.} Що робити, якщо під час медитації виринає старий болісний спогад?

\textbf{В.} Старі спогади, болі, страхи, гнів, образи тощо можуть виринати в медитації. Просто дозвольте їм виринути, не опираючись, не аналізуючи, не оцінюючи й не заперечуючи їх. Лише спостерігайте за ними, не втягуючись. Побачте, що вони не визначають, хто ви. Це згустки несвідомості, що виринають, щоб очиститися у світлі усвідомлення й бути відпущеними з вашої системи. Дозвольте світлу буття визволити страждання.

\textbf{З.} Коли я медитую, часом відчуваю сильний страх. Іноді він мене переповнює, і я не знаю, що робити.

\textbf{В.} Коли в медитації виникає страх, корисно заякорити увагу в чомусь дуже заземлювальному — наприклад, у диханні чи навіть у підошвах ніг. Але не боріться зі страхом, бо це лише його підсилить. Уявіть, що ви Будда під деревом Бодгі або Христос у пустелі: лишайтеся цілком нерухомими й незворушними перед кошмаром тіла-розуму. Він може здаватися дуже реальним, але насправді це не більше ніж переконлива ілюзія. В абсолютній нерухомості без зусиль страх мине сам собою.

\textbf{З.} Що робити, коли під час медитації до мене приходить прозріння чи раптове розуміння якоїсь ситуації?

\textbf{В.} Просто прийміть дароване з вдячністю, нічого не тримаючи. Вірте, що воно буде там, коли вам знадобиться.

\textbf{З.} Я помічаю, що мій розум спонтанно творить образи, майже як сон наяву. Деякі мені подобаються, інші просто випадкові й дратівливі. Що робити?

\textbf{В.} Скеруйте увагу на дихання в животі. Це допоможе не губитися в образах розуму. Тримайте простий намір спочивати в безобразному, мовчазному джерелі, що передує всім образам, думкам та ідеям.

\textbf{З.} Коли я медитую — і навіть коли не медитую — крізь моє тіло проходить багато енергії. Часом це дуже приємно, а часом мене охоплює збудження, і я не сплю ночами. Що це?

\textbf{В.} Не дивно, що в якийсь момент духовного життя ви відчуєте різні форми інтенсивної енергії. Не зачаровуйтеся енергією й не намагайтеся її придушувати чи контролювати, бо це лише схилятиме її до посилення. Укорініть увагу в тому стані, що передує всім формам енергії. Спочивайте в тиші, нерухомості й порожнечі, що передують енергіям тіла й розуму. Укорініть увагу в нижній частині живота. Це допоможе заземлити й інтегрувати енергію.

Також може бути корисно зайнятися дуже заземлювальними справами. Тихі прогулянки на природі, фізичні вправи, масаж підошов ніг тощо. Допоможе будь-що, що відчувається заземлювальним і енергетично заспокійливим. Потрібен час, щоб ваше тіло й нервова система пристосувалися до більшого обсягу енергії, що тече крізь вас. Будьте терплячі. Нервовій системі часто потрібні місяці чи роки, щоб пристосуватися до нового припливу енергії.

\textbf{З.} Іноді я відчуваю таку глибину тиші, у якій усі наміри відпадають, і навіть вказівка дозволяти всьому бути чи спочивати як усвідомлення здається зайвою. Чи можна дозволити всім намірам і технікам відпасти?

\textbf{В.} Часом навіть найтонший намір чи техніка природно відпадають самі, коли медитація сягає певної глибини нерухомості й простоти. Коли ви можете відпустити всі наміри й техніки, не гублячись у розумі й не поринаючи в туманний чи притуплений стан усвідомлення, Справжня медитація відбувається спонтанно. Остаточна форма медитації — це коли той, хто медитує, зникає повністю.

\section*{Самодослідження}

Священного виміру не пізнати словами й ідеями — так само, як не дізнатися, який на смак яблучний пиріг, з'ївши рецепт. Сучасна доба забула, що факти й інформація, хоч які вони корисні, не те саме, що істина чи мудрість, і вже точно не те саме, що безпосередній досвід Дійсності. Ми втратили зв'язок з інтуїтивною мудрістю, народженою з тиші й нерухомості. Тримати питання всередині в мовчазному й терплячому очікуванні — це мистецтво, яке нині опановують рідко. Самодослідження — це міст між его й душею, і далі — до Безкінечного. (Слово «душа» я тут уживаю на позначення сутності, присутності чи буттєвості, якою ви є.)

Самодослідження жодною мірою не антиінтелектуальне й не антираціональне; воно трансраціональне. Тобто воно має силу вивести вас за межі як понятійного розуму, так і зумовленого егоцентричного мислення. Укорінене в нерухомості, самодослідження — це динамічний контрапункт до Справжньої медитації. Медитація м'яка, вона дозволяє віддання, тоді як самодослідження вимагає сміливого й безстрашного запитування. Самодослідження — це спосіб звернутися до найглибших екзистенційних питань, що постають перед кожною людиною: хто я або що я? Що таке життя? Що буде після смерті? Що таке Бог? Що таке абсолютна Істина існування? Або просто: чи знаю я з абсолютною певністю, що ця теперішня думка, вірування, погляд, тлумачення чи судження істинні?

Спільний елемент самодослідження — Істина. Що таке Істина?

Питання про Істину не постає з різних програм его й не стосується їх. Надзвичайно важливо, щоб самодослідження не стало підвладним різним потягам і мотивам его. Глибинні потяги его — почуватися краще й вижити. Але самодослідження належить цілком царині душі — тому виміру буття, народженому з нерухомості й світла, який шукає Істину заради неї самої.

Перший осередок самодослідження — буття. Буття — це ключ, що відмикає царство. Хто я або що я? Поза тілом, розумом, віруванням, заняттям, статтю, роллю, пам'яттю чи історією — що я? Що саме таке «я»?

Відніміть усе, чим «я» не є. Зніміть з «я» всі маски, які воно носить. Що лишилося? Щось? Ніщо? Що усвідомлює це?

У вашому безпосередньому досвіді — щось усвідомлює, чи ніщо усвідомлює? Хтось усвідомлює, чи ніхто не усвідомлює?

Простягайте нитку самодослідження мовчки й терпляче назад крізь усі свої ототожнення, усі вірування про себе, усі потаємні судження й припущення про те, хто ви й що ви. Не поспішайте. Загляньте глибоко в кожне з цих питань. Дозвольте питанням відняти все, чим ви не є. Дозвольте їм розвінчати все, чим ви будь-коли уявляли себе, все, чим ви вважали, що маєте бути, все, чим будь-хто вам казав бути. Простягайте нитку самодослідження назад крізь усі свої уявні тотожності. Ідіть за ниткою назад крізь усе уявлене, усе, за що тримаєтеся, усе, від чого тікаєте. Тоді будьте нерухомі. Спокійте в споглядальній тиші й дозвольте невідомим діянням благодаті звершитися.

Здійснення Істини й Дійсності ніколи не створити розумом; воно завжди приходить як дар благодаті. Самодослідження прибирає хибні сприйняття й ілюзії, роблячи людину доступною для рухів благодаті.

Питання буття відчиняє двері до Дійсності й Істини, але аж ніяк не єдине питання для самодослідження. Питайте все! Не залишайте каменя неперевернутим, жодного припущення недослідженим, жодної форми заперечення незруйнованою.

Досліджуйте кожне питання повільно й навмисно. Вкладіть кожне питання в нерухомість свого буття. Не хапайтеся за швидкі відповіді. Не робіть поспішних висновків. Натомість дозвольте кожному питанню відкрити ваші потаємні вірування й погляди. Дозвольте йому відкрити все, за що ви тримаєтеся й у що вірите, що суперечить тому, що є. Шукайте всі способи, якими прив'язаність до розуму завдає страждання вам та іншим. Приносьте кожне питання, яке ставить розум, у ґрунт нерухомості. Медитуйте над ним, обмірковуйте його; не поспішайте. Не відповідайте на нього розумом. Будьте нерухомі лише з питанням. Будьте дуже, дуже нерухомі.

Сповнені любові до Істини, не дивуйтеся, якщо самодослідження почне поглинати всі ваші потаємні припущення, усі вірування, усі погляди, усі судження, усе, чого ви навчилися з других рук. І не дивуйтеся, якщо більшість ваших духовних ідей теж буде поглинуто, бо саме духовні ідеї найефективніше захищають нас від справді духовного досвіду.

Найбільша поміч — ваша щирість і прагнення Істини понад усе. Вас може знову й знову приголомшувати глибина ілюзії, яку ви знаходите й відкриваєте в собі, але ніколи не затримуйтеся на ній і не судіть себе. Приймайте, прощайте й ідіть далі, бо ваше справжнє буття безкінечне й абсолютне. Воно існує тепер так само, як існувало завжди й існуватиме. Стійте нерухомо у священному полум'ї самодослідження й дозвольте йому відкрити вас до осередку всієї мудрості, народженої з Духу. Виживе лише Істина; усе інше загине.

Сумно, що так мало людей віддають Істині повну міру свого життя. Більшість іде лише до певної межі, а потім задовольняється меншим, ніж цілковите віддання всієї відокремленості. Зрештою ми всі отримуємо те, що найбільше цінуємо, і якщо нам не подобається те, що ми отримали, краще чесно подивитися, що саме ми цінуємо.

Але Істини не бракує ні на мить. Ніколи її не буває більше чи менше, і доступності до неї не буває більше чи менше. Істина в надлишку завжди, за всіх обставин. Вона просто чекає на впізнання. І весь час на її боці.

Питайте свої думки. Питайте свої історії. Питайте свої припущення. Питайте свої погляди. Питайте свої висновки. Питайте їх усі в цілковиту порожнечу, нерухомість і радість. Ключі до свободи у ваших руках. Скористайтеся ними.

\subsection*{Поширені запитання про самодослідження}

\textbf{З.} Самодослідження часто здається мені дуже інтелектуальним, і я просто гублюся в розумі. Чи можна досліджувати так, щоб не губитися в розумі?

\textbf{В.} Так. У самодослідженні є два аспекти, і дуже важливо зрозуміти обидва. Перший аспект — те, що я називаю «кроком назад». Його призначення — відсторонитися від попереднього зумовленого мислення або відняти його. Ви не стільки шукаєте відповідей, скільки відкриваєте й усуваєте попередні зумовлені думки, ідеї та вірування, щоб дати дорогу глибшому здійсненню. Наприклад, через таке спостереження ви можете побачити, що ви не є думками у своєму розумі. Усунувши хибне вірування, ніби якась думка може сказати вам, чим ви є, ви звільняєте простір, щоб глибше розуміння відкрилося саме.

Відкривши й прибравши хибні ідеї розуму, ви готові спочити в нерухомості буття.

У другому аспекті самодослідження ви намагаєтеся дістатися до інтуїтивної ясності й мудрості, що міститься в нерухомості в корені свідомості. Я називаю це цариною благодаті, бо мудрість, яка з неї витікає, завжди отримують як дар, як «ага» чистого розуміння. Відкривши й прибравши хибні ідеї розуму, ви готові спочити в нерухомості буття, не нав'язуючи думки й не шукаючи в думці.

\textbf{З.} Часом самодослідження відчувається дуже живим і важливим для мене, а іноді — механічним, бо серце в ньому не бере участі. Чи мусить самодослідження відбуватися весь час?

\textbf{В.} Щоб самодослідження було справжнім, ви маєте відчувати, що воно життєво вам цікаве. Тож ні, вам не потрібно досліджувати весь час. Це інструмент, який є напохваті щоразу, коли у вас постає життєве питання. Але самодослідження — це також значно більше, ніж техніка; це постава.

Самодослідження — це постава допитливості, що живе у вас, і відображення вашого бажання знати Істину й природу Дійсності. Воно також вимагає свого роду відваги — готовності ставити великі питання, які можуть похитнути самі підвалини вашого життя й поставити вас віч-на-віч із тим, чого ви воліли б уникати. Тож хоч ви й не завжди користуватиметеся технікою самодослідження, життєво важливо жити з поставою допитливості й відваги, що становить серце самодослідження.

\emph{Самодослідження — це мистецтво ставити під сумнів усі свої припущення, вірування й тлумачення, щоб відкрити в розумі простір для виникнення інтуїтивної мудрості. Коли простір відкрито, просто вкладіть питання в нерухомість свідомого буття.}

Спостерігайте. Пильнуйте невідоме вірно. Вирішальні миті прориву приходять тоді, коли ви чекаєте їх найменше.

\section*{Споглядання}

Ми давно забули, що означає споглядати щось. Одним кліком миші ми можемо отримати відповідь — або те, що обіцяє бути відповіддю, — майже на будь-яке питання, яке можемо уявити. Усі найдавніші духовні вчення світу — за один клік або завантаження, і все ж ми лишаємося такими загубленими для себе, такими відрізаними від того, що живить душу, що всі разом задихаємося під вагою свого невігластва й відчуження від священного виміру Життя.

Сучасна доба забула, що факти й інформація, попри всю їхню корисність, не те саме, що мудрість — і вже точно не те саме, що безпосередній досвід Дійсності. Ми втратили зв'язок з інтуїтивною мудрістю, народженою з тиші й нерухомості, і залишилися покинутими в морі інформації, яке не може виконати своєї обіцянки дедалі більшого щастя й повноти.

Споглядання — це мистецтво терпляче тримати слово чи фразу в тиші й нерухомості усвідомлення, доки воно не почне відкривати дедалі глибші значення й розуміння. Споглядання має силу вийти за межі (а не відступити назад від) аналітичного мислення й логіки та відкрити свідомість до такого ладу мудрості й Істини, який можна описати лише як одкровення.

У цьому розділі я вмістив кілька корисних коротких фраз, але предметом споглядання може бути будь-яка частина цієї книги. Візьміть коротку фразу як предмет споглядання й просто тримайте її в усвідомленні якийсь час. Не аналізуйте її й не філософствуйте про неї. І не губіться у власній уяві. Просто тримайте фразу в усвідомленні. Потім будьте нерухомі. Дозвольте її значенню прорости у вас. Потім знову принесіть слово чи фразу в усвідомлення. Тримайте їх там якийсь час, а потім відпустіть і знову будьте нерухомі. З невеликою практикою ви відчуєте, як це робити, і знайдете власний ритм.

Хоч споглядання може здаватися досить простим, воно може бути дуже потужним. У дзенській традиції фрази, питання чи короткі навчальні історії, звані коанами, використовують як предмети споглядання й медитації з великим і потужним наслідком — вони запалюють пробудження, одкровення й просвітлення. Насправді більшість езотеричних духовних систем використовували різні форми споглядання, щоб викликати миті одкровення.

Нижче — фрази для споглядання. Вони починаються з фраз, покликаних запалити певні психологічні прозріння, і ведуть до фраз, що мають привести до дедалі глибших здійснень і одкровень укоріненішої, духовнішої природи. Остання низка фраз стосується абсолютної природи Дійсності, Безкінечного.

\subsection*{Думка й свобода від страждання}

\begin{center}\itshape
Немає такого явища,\\
як абсолютно істинна думка.
\end{center}

Це не означає, що одні думки не істинніші за інші, — лише те, що жодна думка не істинна абсолютно.

\begin{center}\itshape
Те, що є, — це те, що відбувається\\
до того, як у вас з'явилася думка про це.
\end{center}

Зверніть увагу на різницю між тим, що ваш розум думає про цю мить, і цією миттю, якою вона є до будь-якої вашої думки про неї.

\begin{center}\itshape
Страждання виникає тоді, коли ви\\
вірите в думку, що суперечить тому,\\
що є, що було або що може бути.
\end{center}

Переживіть цю мить вільною від тлумачень вашого розуму.

\begin{center}\itshape
Ви — не своя історія.\\
Вони — не ваша історія про них.\\
Світ — не ваша історія про світ.
\end{center}

\begin{center}\itshape
Страждання — це те, як Життя каже вам,\\
що ви опираєтеся або хибно сприймаєте\\
реальне й істинне.
\end{center}

Це спосіб, у який Життя підказує, що ви не в злагоді з тим, що є.

\begin{center}\itshape
Глибше розуміння й прозріння\\
випливають із тихого розуму.
\end{center}

\begin{center}\itshape
Бути щасливим — це\\
жити як невідоме.
\end{center}

\begin{center}\itshape
Усяке справжнє знання виникає з невідомого\\
і є виявом невідомого.
\end{center}

\subsection*{Природа буття}

\begin{center}\itshape
Заглянути всередину й не знайти себе як «я» —\\
це початок того, щоб знайти себе\\
як присутність (буття).
\end{center}

\begin{center}\itshape
Буття (або Дух) універсальне\\
і існує до всіх умов,\\
усіх точок зору, усіх об'єктів свідомості\\
й усіх суб'єктів також.
\end{center}

\begin{center}\itshape
Буття — справжня природа\\
всього.
\end{center}

\begin{center}\itshape
Будучи природою всього,\\
буття не має нічого поза собою.
\end{center}

\begin{center}\itshape
Буття самопізнавальне й усвідомлюване. Просто тепер!
\end{center}

\begin{center}\itshape
Буття не пояснює все;\\
буття — справжня природа всього.
\end{center}

\begin{center}\itshape
Єдине, що здійснює буття, —\\
саме буття.
\end{center}

\begin{center}\itshape
Є лише буття, що живе себе\\
крізь вас, як вас і як усе суще.
\end{center}

\begin{center}\itshape
Буття непороджене й нестворене —\\
джерело й субстанція всього.
\end{center}

\begin{center}\itshape
Буття — наша первісна умова,\\
до всіх егоїстичних витівок, до всякої думки,\\
до всякого опису,\\
до минулого, теперішнього й майбутнього.
\end{center}

\begin{center}\itshape
Та буттєвість, що передує світу простору й часу,\\
є тут, тепер і завжди.\\
Це одна крапля дощу, листок, що падає з дерева,\\
один удар серця.\\
Це світ без світу,\\
субстанція порожнечі.
\end{center}

\begin{center}\itshape
Я Є — це чисте буття.\\
Це останнє визнання Дійсності,\\
що відлунює крізь вічність.
\end{center}

\subsection*{Безкінечне}

\begin{center}\itshape
За его — універсальне буття;\\
за буттям — Безкінечне.
\end{center}

\begin{center}\itshape
Безкінечне — це чистий безформний потенціал,\\
що передує буттю й небуттю, життю й смерті,\\
формі й безформності.
\end{center}

\begin{center}\itshape
Безкінечне — не одне й не багато,\\
не дуалістичне й не недуалістичне,\\
не мирське й не духовне,\\
не своє й не чуже.
\end{center}

\begin{center}\itshape
Безкінечне пізнає себе\\
через простий інтуїтивний погляд, який воно має на себе\\
в кожному аспекті себе.\\
Так воно пізнає себе як цілком непізнаванне\\
й абсолютно присутнє.
\end{center}

\begin{center}\itshape
Здійснити Безкінечне — це втратити свій внутрішній світ.
\end{center}

\begin{center}\itshape
Втратити свій внутрішній світ — це вічна тиша.\\
Це стати сяйвом.
\end{center}

\begin{center}\itshape
Усе добре,\\
і краще, ніж можна уявити.
\end{center}
